\documentclass{llncs}

\usepackage[T1]{fontenc}
\usepackage[utf8]{inputenc}

% Working notes carried over from lit.typ:
%
% == Overview
% - (widely accepted) Runtime verification(properties): [13]
% - Systems: [6, 23]
% - Monitoring Correctness: [28]
% - Fairness of Components: [29]
% - Edge chasing: [45, 54]
% - Simulation relation: [45, 53]
% - Algorithmic idea(seminal): [17, 44]
% - Improvements: [16, 17, 19, 44, 54, 57]
% - Process Calculi: [21, 22, 39-41, 47, 48]
% - Delayed idea: [36, 42]  <- maybe these?
% - Wichtig: [37]
%
% == Todo
% - [17, 44] their deadlock mechanism is inline (cannot satisfy (2)(outline)
%   and (3)transparency)
% - 17deadlock_detection: complete, sound, distributed, Databases, Probes
% - 30deadlock: ??
% - 37formals: Formals
% - 44distributed: Simple implementation using counters
% - 54distributed: Inder, Probes, Edge-chasing
%
% - [43] defines a simulation relation for abstract processes.

\begin{document}

\title{Report: Correct Black-Box Monitors for Distributed
Deadlock Detection}

\author{Sebastian Kl\"ahn}

\institute{Universität Freiburg\\
\email{sk925@email.uni-freiburg.de}}

\maketitle

The Paper `Correct Black-Box Monitors for Distributed Deadlock Detection: Formalisation and Implementation` has two main contributions: Firstly, it provides a novel deadlock algorithm but even more importantly, it provides computer verified formalisation of nodes, networks, monitors and message passing in an RPC framework


\section{The Edge-Chasing Paradigm and Protocol Optimization}
The primary paper's decentralized monitoring logic is directly rooted in
the ``edge-chasing`` paradigm popularized by Chandy, Misra, and Haas
(1983)~\cite{17deadlock_detection}. Chandy et al.\ demonstrated that
deadlocks could be identified without maintaining a centralized global
\emph{wait-for graph} by propagating localized ``probe`` messages across
the network to detect cycles among idle processes. Building upon this,
Mitchell and Merritt~\cite{44distributed} introduced streamlined
algorithms ensuring that only a single process in a cycle detects the
deadlock, thereby vastly simplifying resolution. Crucially, the primary
paper shifts away from forward-moving probes in favor of the
backward-moving probe technique shown by Mitchell and Merritt. This
architectural decision allows to sidestep tricky false positives
that can happen in the forward-moving probe algorithm.

Sinha and Natarajan (1985)~\cite{55priority} make edge-chasing efficient
by incorporating priority-based timestamps. Their algorithm initiates
probes only during ``antagonistic conflicts'' --- when a transaction
waits for one with a lower priority --- drastically reducing
communication overhead and preventing the detection of phantom deadlocks
in fault-free environments. The primary paper leverages these
optimization strategies to ensure its monitoring proxies remain
lightweight and do not become bottlenecks in large-scale RPC systems.

\section{Formal Verification and Operational Semantics}

While early edge-chasing algorithms were conceptually strong, they often
lacked the strict operational semantics required for verifiable
black-box monitoring. To bridge this gap, the primary paper turns to the
formalisms established by Keller (1976)~\cite{37formals}. Keller modeled
parallel computations as \emph{transition systems}, where processes and
computations are represented by discrete state transitions. By adopting
Keller's Labelled Transition System (LTS) semantics, the primary paper
rigorously models networks of Single-threaded Remote Procedure Call
(SRPC) services. This mathematical grounding is what fundamentally
enables the primary paper to mechanically prove monitor transparency and
correctness in the Coq theorem prover.

\section{Mitigating False Positives and System Anomalies}

The demand for mathematical precision in network states is heavily
underscored by Gligor and Shattuck (1980)~\cite{30deadlock}, who exposed
the vulnerabilities of early distributed deadlock protocols. Through
targeted counterexamples, they illustrated how arbitrary network delays
and out-of-order graph updates generate ``ostensibly blocked
transactions,'' ultimately leading to false (phantom) deadlocks. Gligor
and Shattuck's cautionary findings directly inform the primary paper's
stringent correctness criteria. By demanding both ``soundness'' and
``completeness,'' the primary paper establishes verifiable invariants
ensuring that monitors only report actual deadlocks, neutralizing the
anomalies Gligor and Shattuck originally identified.

\section{Transactional Dependencies and Concurrency Control}

Finally, managing distributed deadlocks requires a precise understanding
of the underlying resource dependencies, a concept comprehensively
addressed by Moss (1981)~\cite{54distributed}. Moss extended traditional
single-level transaction systems to support nested transactions,
introducing strict two-phase locking rules and state restoration
algorithms to maintain failure atomicity and concurrency control.
Although the primary paper is tailored toward RPC service deadlocks
rather than database transaction management, Moss's framework for
handling distributed locking and wait dependencies clarifies the complex
hierarchical network behaviors that modern edge-chasing monitors must
correctly interpret and navigate.

\bibliographystyle{splncs04}
\bibliography{literature/all}

\end{document}
